\subsubsection{Geometric domain}
\label{subsec:geo-domain}
\newcommand{\Caddy}{\textit{Caddy}\xspace}
\newcommand{\CoreCaddy}{\textit{Core Caddy}\xspace}

\sz~\citep{szalinski} is an \eqsat-driven
  tool that shrinks 3D CAD (Computer-Aided Design) programs
  by performing rewrites
  over a language called \Caddy.
  \Caddy expresses CAD programs with
  primitives for
  constructive solid geometry (e.g., \C{Cube}, \C{Scale}, \C{Union}),
  basic rational arithmetic,
  various list constructors,
  and inverse transformations.
\sz shrinks \Caddy programs using
  two rulesets:
  a set of CAD
  identities
  and a set of custom procedural rewrite rules
  which discover opportunities to use inverse
  transformations.

We focus on synthesizing
  the CAD identities,
  which help \sz expose
  hidden structure in its input programs.
Synthesizing these identities using
  traditional rule inference algorithms
  requires a full CAD interpreter,
  which is difficult to implement and
  too slow for use in rule synthesis.
Instead, we leverage \enumo's \ff approach.

\paragraph*{Synthesizing CAD identities.}
Solid geometry can be
  represented mathematically via
  a function representation (F-Rep).
An F-Rep is a function $f(x, y, z)$
  which interprets an arithmetic expression
  over $x$, $y$, and $z$ as the geometric solid
  defined where $f$ is positive~\citep{frep}.
For example, the unit sphere can be represented
  by $1 - x^2 - y^2 - z^2$.

We use \enumo to synthesize a set of CAD
  identities by
  fast-forwarding from a small set of 5
  translational rules from CAD to F-Rep
  combined with 15 rules over the F-Rep domain.
Our prior rules for F-Rep included 6 rules
  over rationals and 9 substitution rules
  over operators that allowed it to express \Caddy
  transformations.
Here are two examples of these rules:
\[
  \mathit{Sphere}(r) \rewritesto  1 -  \left(\frac{x}{r}\right)^2 - \left(\frac{y}{r}\right)^2 - \left(\frac{z}{r}\right)^2
  \qquad
  \mathit{Scale}([w, h, l], e) \rewritesto e\left[ x \mapsto \frac{x}{w}\right]\left[ y \mapsto \frac{y}{h}\right]\left[ z \mapsto \frac{z}{l}\right]
\]
Using \enumo's guided search and fast-forwarding,
  we synthesized a set of 10 large bidirectional
  CAD identities of up to 13 atoms:
\begin{enumerate}
  \item \foottt{(Trans (Vec3 0 0 0) a) \rewritesboth \xspace a}
  \item \foottt{(Scale (Vec3 1 1 1) a) \rewritesboth \xspace a}
  \item \foottt{(Cylinder (Vec3 b b b) a true) \rewritesboth \linebreak
    (Scale (Vec3 b b b) (Cylinder (Vec3 1 1 1) a true))}
  \item \foottt{(Sphere b a) \rewritesboth \xspace (Scale (Vec3 b b b) (Sphere 1 a))}
  \item \foottt{(Scale (Vec3 f e d) (Cube (Vec3 c b a) false)) \rewritesboth \linebreak
    (Scale (Vec3 (* f c) (* e b) (* d a)) (Cube (Vec3 1 1 1) false))}
  \item \foottt{(Cube (Vec3 (* f e) (* d c) (* b a)) false) \rewritesboth \linebreak
    (Scale (Vec3 f d b) (Cube (Vec3 e c a) false))}
  \item \foottt{(Trans (Vec3 g f e) (Scale (Vec3 d c b) a)) \rewritesboth \linebreak
    (Scale (Vec3 d c b) (Trans (Vec3 (/ g d) (/ f c) (/ e b)) a))}
  \item \foottt{(Scale (Vec3 g f e) (Trans (Vec3 d c b) a)) \rewritesboth \linebreak
    (Trans (Vec3 (* d g) (* c f) (* b e)) (Scale (Vec3 g f e) a))}
  \item \foottt{(Scale (Vec3 g f e) (Scale (Vec3 d c b) a)) \rewritesboth \linebreak
    (Scale (Vec3 (* g d) (* f c) (* e b)) a)}
  \item \foottt{(Trans (Vec3 (+ g f) (+ e d) (+ c b)) a) \rewritesboth \linebreak
    (Trans (Vec3 g e c) (Trans (Vec3 f d b) a))}
\end{enumerate}

The workload we used is shown in \autoref{fig:szalinski-workload}.

\paragraph*{Results.}
We evaluate \sz's performance
  on a set of
  benchmarks taken from the paper (Table 2 of \citealt{szalinski})\footnote{
    \citet{szalinski} states that
     these benchmarks were decompiled using the
     Reincarnate~\citep{reincarnate} mesh decompiler.}.
The results are shown in \autoref{fig:szalinski-results}.
Using \enumo's synthesized rules for CAD identities, \sz is able
  to shrink input benchmarks by 87\% on average, while the original, handwritten
  rules shrink input benchmarks by 90\% on average.
Upon closer inspection, we found that a subset of the \enumo-synthesized
  rules matches the performance of \sz's handwritten CAD identities.
There are some rules in the \enumo ruleset that
  cannot be derived from \sz's CAD identities.
However, we posit that these additional rules
  are not very useful for the \sz benchmark tests, so their
  presence in the \enumo ruleset leads to worse
  performance at low node limits.

\begin{figure}

  \begin{lstlisting}[
    name=szalinski-workload-code,
    language=Python,
    basicstyle=\footnotesize\ttfamily,
    numbers=left,
    xleftmargin=2.5em]
  def iter_szalinski(n):
    lang = { (AFFINE VEC SOLID) (Cube VEC) (Cylinder VEC) (Sphere SCALAR) }
    return iter_metric(lang, "SOLID", Depth, n)
             .plug("VEC", { (Vec3 0 0 0) (Vec3 1 1 1) (Vec3 a a a) (Vec3 a b c)
                           (Vec3 d e f) (Vec3 (BOP a d) (BOP b e) (BOP c f)) })
             .plug("AFFINE", { Scale Trans })
             .plug("SCALAR", { a 1 })
             .plug("BOP", { + * / })

  cad_idents = []
  for i in [2, 3]:
      wkld = iter_szalinski(i)
      chosen = fast_forward(wkld, frep_rules, cad_idents)
      cad_idents.extend(chosen)
  \end{lstlisting}
  \caption{
    An \enumo program for learning CAD identities. \texttt{iter_szalinski}
    is a function that constructs workloads for \Caddy expressions.
  }
  \label{fig:szalinski-workload}
  \label{fig:szalinski-recipe}
\end{figure}

\begin{table}[h]
    \footnotesize
  \begin{tabular}{lccc}
    Program Id & No Identities & \sz's Identities & \enumo's Identities \\
\midrule
TackleBox & 79\% (60) & 91\% (26) & 85\% (41) \\
SDCardRack & 72\% (57) & 87\% (26) & 86\% (28) \\
SingleRowHolder & 84\% (31) & 92\% (16) & \textbf{92\% (16)} \\
CircleCell & 61\% (31) & 80\% (16) & \textbf{80\% (16)} \\
CNCBitCase & 88\% (27) & 93\% (15) & 93\% (16) \\
CassetteStorage & 81\% (27) & 89\% (15) & \textbf{89\% (15)} \\
RaspberryPiCover & 90\% (27) & 96\% (12) & \textbf{96\% (12)} \\
ChargingStation & 81\% (27) & 89\% (15) & 82\% (25) \\
CardFramer & 52\% (83) & 76\% (42) & \textbf{76\% (42)} \\
HexWrenchHolder & 90\% (31) & 95\% (16) & 84\% (52) \\
\midrule
Average & 80\% (40.1) & 90\% (19.9) & 87\% (26.3) \\
    \end{tabular}

\caption{End-to-end evaluation of
  \sz{} (Table 2 of \citealt{szalinski})
  using CAD identities from \enumo.
The inputs are flat CAD programs from  Reincarnate~\citep{reincarnate}.
For each set of identities,
  we show the percentages by which the initial program's AST sizes decreased
  (higher is better), and in parentheses, the output AST sizes (lower is better).
\texttt{No Identities}, \texttt{\sz's Identities}, and \texttt{\enumo's Identities} show the results of running \sz
  without CAD identities, with all identities enabled, and with
  \enumo's synthesized CAD identities in place of the original ones, respectively.
  \enumo's CAD identities exactly matched the
  performance of Szalinski's on 5 of 10 inputs.}

\label{fig:szalinski-results}
\end{table}

\autoref{fig:szalinski-results} shows that
  \enumo's CAD identities are able to closely match the performance
  of the handwritten CAD identities.
Fast-forwarding is effective at synthesizing rules for
  constructive geometry,
  which is a domain containing rewrite rules with term sizes that
  cannot be exhaustively enumerated,
  for which writing an interpreter is challenging,
  and which prior work did not support.
